由 [SGA 1, XIII 2.8], 存在特化同态 (依赖于 $T$ 中从 $\bar\eta$ 到 $o$ 的一条道路, 因而只确定到共轭)
\begin{equation}\tag{6.1.3}
\pi_1^{\mathrm{mod}}(D_{\bar\eta}-D_{\bar\eta}\cap Z,\bar a)
\xrightarrow{\ \bar\eta\longrightarrow o\ }
\pi_1^{\mathrm{mod}}(D_o-D_o\cap Z,\bar a)
\end{equation}
使得下列断言成立.

\subsubsection*{6.1.4.}
局部单值同态 (6.1.2.1)
\[
\prod_{\ell\ne\operatorname{char}(k)}\mathbb Z_\ell(1)
\longrightarrow
\pi_1^{\mathrm{mod}}(D_o-D_o\cap Z,\bar a)
\]
由 $D_{\bar\eta}$ 的局部单值同态与特化箭头 (6.1.3) 复合得到:
\begin{equation}\tag{6.1.4.1}
\prod_{\ell\ne\operatorname{char}(k)}\mathbb Z_\ell(1)
\longrightarrow
\pi_1^{\mathrm{mod}}(D_{\bar\eta}-D_{\bar\eta}\cap Z,\bar a)
\xrightarrow{\ \bar\eta\longrightarrow o\ }
\pi_1^{\mathrm{mod}}(D_o-D_o\cap Z,\bar a).
\end{equation}

\subsubsection*{6.1.5.}
下图交换:
\begin{equation}\tag{6.1.5}
\begin{tikzcd}[column sep=huge,row sep=huge]
{} & \pi_1^{\mathrm{mod}}(\mathbb P^r-Z,\bar a) & {} \\
\pi_1^{\mathrm{mod}}(D_\eta-D_\eta\cap Z,\bar a)
  \arrow[ur,"i_{\eta*}"]
& {}
& \pi_1^{\mathrm{mod}}(D_o-D_o\cap Z,\bar a)
  \arrow[ul,"i_{o*}"'] \\
{} & \pi_1^{\mathrm{mod}}(D_{\bar\eta}-D_{\bar\eta}\cap Z,\bar a)
  \arrow[uu,"i_{\bar\eta*}"']
  \arrow[ul]
  \arrow[ur,"\bar\eta\longrightarrow o"'] & {} .
\end{tikzcd}
\end{equation}

暂且承认这一点, 先证明 6.1.1. 由 6.1.5 的交换性, 只须证明
\[
i_{\bar\eta*}:
\pi_1^{\mathrm{mod}}(D_{\bar\eta}-D_{\bar\eta}\cap Z,\bar a)
\longrightarrow
\pi_1^{\mathrm{mod}}(\mathbb P^r-Z,\bar a)
\]
为满射. 为此, 只须证明\underline{通常}基本群的相应同态
\begin{equation}\tag{6.1.6}
i_{\bar\eta*}:
\pi_1(D_{\bar\eta}-D_{\bar\eta}\cap Z,\bar a)
\longrightarrow
\pi_1(\mathbb P^r-Z,\bar a)
\end{equation}
为满射. 换言之, 对每个 $E$ 连通的平展覆盖 $E\to\mathbb P^r-Z$ (因 $\mathbb P^r-Z$ 光滑, 从而单分支 [EGA IV, 18.10.3], 故也不可约), 它的逆像
\[
E_{\bar\eta}
=E\times_{\mathbb P^r-Z}(D_{\bar\eta}-D_{\bar\eta}\cap Z)
\]
在 $D_{\bar\eta}-D_{\bar\eta}\cap Z$ 上须保持几何不可约. 只须将以下 Bertini 定理的变体 ([EGA V, 原文未填写引用位置]) 应用于复合 $E\to\mathbb P^r-Z\lhook\joinrel\longrightarrow\mathbb P^r$ 及 $s=1$.

\subsubsection*{6.1.6.1. Bertini 定理.}
设 $k$ 为域, $\mathbb P^r=\mathbb P^r_k$, $f:X\to\mathbb P^r$ 为非常值态射, 其中 $X$ 为几何不可约代数概形. 对每个整数 $s<\dim\overline{f(X)}$, 以 $L_t$ 表示 $\mathbb P^r$ 中余维为 $s$ 的\underline{一般}线性子空间. 则 $X_t=f^{-1}(L_t)$ 几何不可约.
